% AUTO-GENERATED by scripts/gen_blueprint.py (docs/all-axioms.txt + docs/axiom-report.txt + curated metadata).
% Do not edit by hand; re-run the generator.

\chapter{Concrete genus theorems}

Explicit projective-curve models on which the genus is computed directly. The first three are fully axiom-free; the even hyperelliptic case rests on the Liouville L2/L3 frontier.

\begin{theorem}[$\operatorname{genus}\PP^1 = 0$]
  \label{thm:genus-P1}
  \lean{Jacobians.ProjectiveCurve.genus_projectiveLine_eq_zero}
  \leanok
  Proved in Lean.
\end{theorem}

\begin{theorem}[$\Omega^1(\PP^1) = 0$]
  \label{thm:forms-P1}
  \lean{Jacobians.ProjectiveCurve.HolomorphicOneForm_projectiveLine_eq_zero}
  \leanok
  Proved in Lean.
\end{theorem}

\begin{theorem}[genus of an elliptic curve $= 1$]
  \label{thm:genus-E}
  \lean{Jacobians.ProjectiveCurve.genus_Elliptic_eq_one}
  \leanok
  Proved in Lean.
\end{theorem}

\begin{theorem}[genus of an even hyperelliptic curve $= N/2-1$]
  \label{thm:genus-hyp}
  \lean{Jacobians.Extensions.HyperellipticEven.genus_HyperellipticEven_eq}
  \leanok
  \uses{ax:AX-HyperellipticAffine-connected,ax:AX-HyperellipticOneForm-eq-form,ax:affineLiftChart-compat-infinityLiftChart,ax:contDiffOn-symm-toOpenPartialHomeomorph,ax:infinityLiftChart-compat-affineLiftChart,ax:polynomialLocalHomeomorph-no-critical-in-source,ax:squareLocalHomeomorph-zero-notMem-source}
  Proved in Lean, resting on the axiom frontier below.
\end{theorem}

\chapter{The abstract Jacobian (Buzzard's challenge)}

Buzzard's target declarations, proved in Lean modulo the named axioms below --- the classical Riemann-surface theory each rests on.

\begin{definition}[the genus $\dim_{\CC} \Omega^1(X)$]
  \label{def:genus}
  \lean{genus}
  \leanok
  Proved in Lean.
\end{definition}

\begin{definition}[the Jacobian $\Omega^1(X)^* / H_1(X;\mathbb{Z})$]
  \label{def:jac}
  \lean{Jacobian}
  \leanok
  \uses{ax:loopIntegralToH1}
  Proved in Lean, resting on the axiom frontier below.
\end{definition}

\begin{definition}[the Abel--Jacobi map $X \to \operatorname{Jac}(X)$]
  \label{def:aj}
  \lean{Jacobian.ofCurve}
  \leanok
  \uses{ax:loopIntegralToH1,ax:pathIntegralBasepointFunctional}
  Proved in Lean, resting on the axiom frontier below.
\end{definition}

\begin{definition}[the degree of a holomorphic map]
  \label{def:deg}
  \lean{ContMDiff.degree}
  \leanok
  \uses{ax:AX-BranchLocus,ax:localOrder}
  Proved in Lean, resting on the axiom frontier below.
\end{definition}

\begin{theorem}[genus $0 \iff$ homeomorphic to $\PP^1$]
  \label{thm:g0}
  \lean{genus_eq_zero_iff_homeo}
  \leanok
  \uses{ax:AX-genus-eq-zero-iff-homeo}
  Proved in Lean, resting on the axiom frontier below.
\end{theorem}

\begin{theorem}[Abel--Jacobi sends the basepoint to $0$]
  \label{thm:aj-self}
  \lean{Jacobian.ofCurve_self}
  \leanok
  \uses{ax:loopIntegralToH1,ax:pathIntegralBasepointFunctional}
  Proved in Lean, resting on the axiom frontier below.
\end{theorem}

\begin{theorem}[Abel--Jacobi is injective (genus $>0$)]
  \label{thm:aj-inj}
  \lean{Jacobian.ofCurve_inj}
  \leanok
  \uses{ax:AX-ofCurve-inj,ax:loopIntegralToH1,ax:pathIntegralBasepointFunctional}
  Proved in Lean, resting on the axiom frontier below.
\end{theorem}

\begin{theorem}[Abel--Jacobi is holomorphic]
  \label{thm:aj-smooth}
  \lean{Jacobian.ofCurve_contMDiff}
  \leanok
  \uses{ax:AX-PeriodLattice,ax:AX-ofCurve-contMDiff,ax:instPeriodLatticeDiscrete,ax:loopIntegralToH1,ax:pathIntegralBasepointFunctional}
  Proved in Lean, resting on the axiom frontier below.
\end{theorem}

\begin{definition}[pushforward of $1$-forms / points]
  \label{def:push}
  \lean{Jacobian.pushforward}
  \leanok
  \uses{ax:AX-pushforwardAmbient-preserves-lattice,ax:loopIntegralToH1,ax:pullbackOneForm}
  Proved in Lean, resting on the axiom frontier below.
\end{definition}

\begin{definition}[pullback of $1$-forms / points]
  \label{def:pull}
  \lean{Jacobian.pullback}
  \leanok
  \uses{ax:AX-pullbackAmbient-preserves-lattice,ax:loopIntegralToH1,ax:pushforwardOneForm}
  Proved in Lean, resting on the axiom frontier below.
\end{definition}

\begin{theorem}[pushforward is holomorphic]
  \label{thm:push-smooth}
  \lean{Jacobian.pushforward_contMDiff}
  \leanok
  \uses{ax:AX-PeriodLattice,ax:AX-pushforward-contMDiff,ax:AX-pushforwardAmbient-preserves-lattice,ax:instPeriodLatticeDiscrete,ax:loopIntegralToH1,ax:pullbackOneForm}
  Proved in Lean, resting on the axiom frontier below.
\end{theorem}

\begin{theorem}[pullback is holomorphic]
  \label{thm:pull-smooth}
  \lean{Jacobian.pullback_contMDiff}
  \leanok
  \uses{ax:AX-PeriodLattice,ax:AX-pullback-contMDiff,ax:AX-pullbackAmbient-preserves-lattice,ax:instPeriodLatticeDiscrete,ax:loopIntegralToH1,ax:pushforwardOneForm}
  Proved in Lean, resting on the axiom frontier below.
\end{theorem}

\begin{theorem}[pushforward of the identity]
  \label{thm:push-id}
  \lean{Jacobian.pushforward_id_apply}
  \leanok
  \uses{ax:AX-pullbackOneForm-id,ax:AX-pushforwardAmbient-preserves-lattice,ax:loopIntegralToH1,ax:pullbackOneForm}
  Proved in Lean, resting on the axiom frontier below.
\end{theorem}

\begin{theorem}[pushforward is functorial]
  \label{thm:push-comp}
  \lean{Jacobian.pushforward_comp_apply}
  \leanok
  \uses{ax:AX-pullbackOneForm-comp,ax:AX-pushforwardAmbient-preserves-lattice,ax:loopIntegralToH1,ax:pullbackOneForm}
  Proved in Lean, resting on the axiom frontier below.
\end{theorem}

\begin{theorem}[pullback of the identity]
  \label{thm:pull-id}
  \lean{Jacobian.pullback_id_apply}
  \leanok
  \uses{ax:AX-pullbackAmbient-preserves-lattice,ax:AX-pushforwardOneForm-id,ax:loopIntegralToH1,ax:pushforwardOneForm}
  Proved in Lean, resting on the axiom frontier below.
\end{theorem}

\begin{theorem}[pullback is functorial]
  \label{thm:pull-comp}
  \lean{Jacobian.pullback_comp_apply}
  \leanok
  \uses{ax:AX-pullbackAmbient-preserves-lattice,ax:AX-pushforwardOneForm-comp,ax:loopIntegralToH1,ax:pushforwardOneForm}
  Proved in Lean, resting on the axiom frontier below.
\end{theorem}

\begin{theorem}[$\text{push}\circ\text{pull} = \deg$]
  \label{thm:push-pull}
  \lean{Jacobian.pushforward_pullback}
  \leanok
  \uses{ax:AX-BranchLocus,ax:AX-pullbackAmbient-preserves-lattice,ax:AX-pushforward-pullback,ax:AX-pushforwardAmbient-preserves-lattice,ax:localOrder,ax:loopIntegralToH1,ax:pullbackOneForm,ax:pushforwardOneForm}
  Proved in Lean, resting on the axiom frontier below.
\end{theorem}

\chapter{The axiom frontier}

Every node below is currently an \texttt{axiom} --- a staging point to discharge into a Lean proof (a subproject). Edges into these nodes are machine-checked. Full audit: \texttt{AXIOM\_AUDIT.md}.

\section{Class 1 --- textbook-standard (trusted)}

Standard classical theorems with citations; discharge $=$ port the proof or import from Mathlib.

\begin{lemma}[AX\_AbelTheorem]
  \label{ax:AX-AbelTheorem}
  \lean{Jacobians.Axioms.AX_AbelTheorem}
  Abel's theorem (Forster \S21).
\end{lemma}

\begin{lemma}[AX\_AnalyticCycleBasis]
  \label{ax:AX-AnalyticCycleBasis}
  \lean{Jacobians.Axioms.AX_AnalyticCycleBasis}
  symplectic $H_1$ basis (standard).
\end{lemma}

\begin{lemma}[AX\_BranchLocus]
  \label{ax:AX-BranchLocus}
  \lean{Jacobians.Axioms.AX_BranchLocus}
  branch-locus / degree finiteness (Miranda Ch.~II).
\end{lemma}

\begin{lemma}[AX\_IntersectionForm\_alternating]
  \label{ax:AX-IntersectionForm-alternating}
  \lean{Jacobians.Axioms.AX_IntersectionForm_alternating}
  cup product on $H_1$ is alternating.
\end{lemma}

\begin{lemma}[AX\_IntersectionForm\_perfect]
  \label{ax:AX-IntersectionForm-perfect}
  \lean{Jacobians.Axioms.AX_IntersectionForm_perfect}
  Poincar\'e duality / unimodularity.
\end{lemma}

\begin{lemma}[AX\_PeriodLattice]
  \label{ax:AX-PeriodLattice}
  \lean{Jacobians.Axioms.AX_PeriodLattice}
  the period lattice is a full $\mathbb{Z}$-lattice.
\end{lemma}

\begin{lemma}[AX\_PluckerFormula]
  \label{ax:AX-PluckerFormula}
  \lean{Jacobians.Axioms.AX_PluckerFormula}
  Pl\"ucker formula (Griffiths--Harris Ch.~2).
\end{lemma}

\begin{lemma}[AX\_RiemannBilinear]
  \label{ax:AX-RiemannBilinear}
  \lean{Jacobians.Axioms.AX_RiemannBilinear}
  Riemann bilinear relations (Griffiths--Harris Ch.~2).
\end{lemma}

\begin{lemma}[AX\_RiemannRoch]
  \label{ax:AX-RiemannRoch}
  \lean{Jacobians.Axioms.AX_RiemannRoch}
  Riemann--Roch (Forster \S16).
\end{lemma}

\begin{lemma}[AX\_SerreDuality]
  \label{ax:AX-SerreDuality}
  \lean{Jacobians.Axioms.AX_SerreDuality}
  Serre duality (Forster \S17).
\end{lemma}

\begin{lemma}[AX\_genus\_eq\_zero\_iff\_homeo]
  \label{ax:AX-genus-eq-zero-iff-homeo}
  \lean{Jacobians.Axioms.AX_genus_eq_zero_iff_homeo}
  uniformization, genus $0$ (Forster \S27).
\end{lemma}

\begin{lemma}[instPeriodLatticeDiscrete]
  \label{ax:instPeriodLatticeDiscrete}
  \lean{Jacobians.Axioms.instPeriodLatticeDiscrete}
  discreteness of the period lattice.
\end{lemma}

\begin{lemma}[ambientPhi\_ambientPsi\_eq]
  \label{ax:ambientPhi-ambientPsi-eq}
  \lean{Jacobians.Vendor.Kirov.ambientPhi_ambientPsi_eq}
  degree identity (Kirov handoff).
\end{lemma}

\begin{lemma}[genus\_eq\_zero\_iff\_homeo]
  \label{ax:genus-eq-zero-iff-homeo}
  \lean{Jacobians.Vendor.Kirov.genus_eq_zero_iff_homeo}
  uniformization (Kirov handoff).
\end{lemma}

\section{Class 2a --- data-existence (spec under review)}

``This object exists with spec $S$.'' The spec, not the proof, is the question.

\begin{lemma}[AX\_pathIntegral\_local\_antiderivative]
  \label{ax:AX-pathIntegral-local-antiderivative}
  \lean{Jacobians.Axioms.AX_pathIntegral_local_antiderivative}
  chart-local FTC binding the functional to the cocycle.
\end{lemma}

\begin{lemma}[abelJacobiDiv]
  \label{ax:abelJacobiDiv}
  \lean{Jacobians.Axioms.abelJacobiDiv}
  divisor-level Abel--Jacobi.
\end{lemma}

\begin{lemma}[intersectionForm]
  \label{ax:intersectionForm}
  \lean{Jacobians.Axioms.intersectionForm}
  the $H_1$ intersection pairing.
\end{lemma}

\begin{lemma}[localOrder]
  \label{ax:localOrder}
  \lean{Jacobians.Axioms.localOrder}
  local multiplicity of a holomorphic map.
\end{lemma}

\begin{lemma}[pathIntegralBasepointFunctional]
  \label{ax:pathIntegralBasepointFunctional}
  \lean{Jacobians.Axioms.pathIntegralBasepointFunctional}
  the path-integral functional; \emph{opaque}.
\end{lemma}

\begin{lemma}[pullbackOneForm]
  \label{ax:pullbackOneForm}
  \lean{Jacobians.Axioms.pullbackOneForm}
  pullback of holomorphic $1$-forms.
\end{lemma}

\begin{lemma}[pushforwardOneForm]
  \label{ax:pushforwardOneForm}
  \lean{Jacobians.Axioms.pushforwardOneForm}
  trace (pushforward) of holomorphic $1$-forms.
\end{lemma}

\begin{lemma}[loopIntegralToH1]
  \label{ax:loopIntegralToH1}
  \lean{Jacobians.RiemannSurface.loopIntegralToH1}
  $H_1$-level period descent.
\end{lemma}

\begin{lemma}[Sheaf-cohomology type stubs (13)]
  \label{clu:sheaf}
  \lean{Jacobians.Axioms.Divisor}
  The line-bundle / divisor / sheaf-cohomology layer ($\mathrm{Divisor}$, $\mathrm{LineBundle}$, $H^0$, $H^1$, $\ldots$) is axiomatized as types + instances --- the part of the encoding most in question. \emph{Cluster of 13 axioms.}
\end{lemma}

\begin{lemma}[Kirov line-integral bridge path (6)]
  \label{clu:bridge}
  \lean{Jacobians.Bridge.bridgePath}
  Path-selection data ($\mathrm{bridgePath}$ + continuity / differentiability / endpoints / integrability) for the Kirov line-integral bridge. \emph{Cluster of 6 axioms.}
\end{lemma}

\section{Class 2b --- definition-asserting (may mask a bad definition)}

``My construction behaves correctly.'' Validate on a concrete witness.

\begin{lemma}[AX\_ofCurve\_contMDiff]
  \label{ax:AX-ofCurve-contMDiff}
  \lean{Jacobians.Axioms.AX_ofCurve_contMDiff}
  Abel--Jacobi smoothness.
\end{lemma}

\begin{lemma}[AX\_ofCurve\_inj]
  \label{ax:AX-ofCurve-inj}
  \lean{Jacobians.Axioms.AX_ofCurve_inj}
  Abel--Jacobi injectivity; \emph{opaque-blocked}.
\end{lemma}

\begin{lemma}[AX\_pullbackAmbient\_preserves\_lattice]
  \label{ax:AX-pullbackAmbient-preserves-lattice}
  \lean{Jacobians.Axioms.AX_pullbackAmbient_preserves_lattice}
  period-map naturality (pullback).
\end{lemma}

\begin{lemma}[AX\_pullbackOneForm\_comp]
  \label{ax:AX-pullbackOneForm-comp}
  \lean{Jacobians.Axioms.AX_pullbackOneForm_comp}
  pullback is contravariant.
\end{lemma}

\begin{lemma}[AX\_pullbackOneForm\_id]
  \label{ax:AX-pullbackOneForm-id}
  \lean{Jacobians.Axioms.AX_pullbackOneForm_id}
  pullback preserves identity.
\end{lemma}

\begin{lemma}[AX\_pullback\_contMDiff]
  \label{ax:AX-pullback-contMDiff}
  \lean{Jacobians.Axioms.AX_pullback_contMDiff}
  pullback smoothness.
\end{lemma}

\begin{lemma}[AX\_pushforwardAmbient\_preserves\_lattice]
  \label{ax:AX-pushforwardAmbient-preserves-lattice}
  \lean{Jacobians.Axioms.AX_pushforwardAmbient_preserves_lattice}
  period-map naturality (pushforward).
\end{lemma}

\begin{lemma}[AX\_pushforwardOneForm\_comp]
  \label{ax:AX-pushforwardOneForm-comp}
  \lean{Jacobians.Axioms.AX_pushforwardOneForm_comp}
  trace is covariant.
\end{lemma}

\begin{lemma}[AX\_pushforwardOneForm\_id]
  \label{ax:AX-pushforwardOneForm-id}
  \lean{Jacobians.Axioms.AX_pushforwardOneForm_id}
  trace preserves identity.
\end{lemma}

\begin{lemma}[AX\_pushforward\_contMDiff]
  \label{ax:AX-pushforward-contMDiff}
  \lean{Jacobians.Axioms.AX_pushforward_contMDiff}
  pushforward smoothness.
\end{lemma}

\begin{lemma}[AX\_pushforward\_pullback]
  \label{ax:AX-pushforward-pullback}
  \lean{Jacobians.Axioms.AX_pushforward_pullback}
  push $\circ$ pull $=$ degree.
\end{lemma}

\section{Class 2c --- atlas / structure (curve-specific chart work)}

Real chart / manifold constructions; classically true, discharge is substantial chart calculus. The prime subprojects.

\begin{lemma}[contDiffOn\_symm\_toOpenPartialHomeomorph]
  \label{ax:contDiffOn-symm-toOpenPartialHomeomorph}
  \lean{Jacobians.GeneralResults.contDiffOn_symm_toOpenPartialHomeomorph}
  narrow inverse-function-theorem gap.
\end{lemma}

\begin{lemma}[AX\_H1\_ProjectiveLine\_trivial]
  \label{ax:AX-H1-ProjectiveLine-trivial}
  \lean{Jacobians.ProjectiveCurve.AX_H1_ProjectiveLine_trivial}
  $H_1(\PP^1)$ is trivial.
\end{lemma}

\begin{lemma}[AX\_HyperellipticAffine\_connected]
  \label{ax:AX-HyperellipticAffine-connected}
  \lean{Jacobians.ProjectiveCurve.HyperellipticAffine.AX_HyperellipticAffine_connected}
  the affine hyperelliptic curve is connected.
\end{lemma}

\begin{lemma}[polynomialLocalHomeomorph\_no\_critical\_in\_source]
  \label{ax:polynomialLocalHomeomorph-no-critical-in-source}
  \lean{Jacobians.ProjectiveCurve.HyperellipticAffine.polynomialLocalHomeomorph_no_critical_in_source}
  affine-form IFT shape (polynomial branch).
\end{lemma}

\begin{lemma}[squareLocalHomeomorph\_zero\_notMem\_source]
  \label{ax:squareLocalHomeomorph-zero-notMem-source}
  \lean{Jacobians.ProjectiveCurve.HyperellipticAffine.squareLocalHomeomorph_zero_notMem_source}
  affine-form IFT shape ($\sqrt{f}$ branch).
\end{lemma}

\begin{lemma}[affineLiftChart\_compat\_infinityLiftChart]
  \label{ax:affineLiftChart-compat-infinityLiftChart}
  \lean{Jacobians.ProjectiveCurve.HyperellipticEvenProj.affineLiftChart_compat_infinityLiftChart}
  cross-summand chart compat, affine $\to\infty$ (SP-1).
\end{lemma}

\begin{lemma}[infinityLiftChart\_compat\_affineLiftChart]
  \label{ax:infinityLiftChart-compat-affineLiftChart}
  \lean{Jacobians.ProjectiveCurve.HyperellipticEvenProj.infinityLiftChart_compat_affineLiftChart}
  cross-summand chart compat, $\infty\to$ affine (SP-2).
\end{lemma}

\begin{lemma}[Unified hyperelliptic curve: type + instances (11)]
  \label{clu:hyp}
  \lean{Jacobians.ProjectiveCurve.AX_Hyperelliptic_evenEquiv}
  The abstract $\mathrm{Hyperelliptic}$ type, its 7 Riemann-surface typeclass instances, and $\mathrm{oddEquiv}/\mathrm{evenEquiv}/\mathrm{genus}$. \emph{Cluster of 11 axioms.}
\end{lemma}

\begin{lemma}[Plane curve: type + instances (11)]
  \label{clu:plane}
  \lean{Jacobians.ProjectiveCurve.PlaneCurve}
  The $\mathrm{PlaneCurve}$ type, its 7 instances, and the affine connected/noncompact/nonempty facts. \emph{Cluster of 11 axioms.}
\end{lemma}

\begin{lemma}[Odd-atlas infinity chart (7)]
  \label{clu:odd}
  \lean{Jacobians.ProjectiveCurve.HyperellipticOdd.affineLiftProjX_compat_infinityChart}
  The odd-degree atlas's chart at infinity ($\mathrm{infinityChart}$, its inverse, four compatibility facts, and membership). \emph{Cluster of 7 axioms.}
\end{lemma}

\begin{lemma}[Elliptic period witnesses (3)]
  \label{clu:ell}
  \lean{Jacobians.ProjectiveCurve.AX_Elliptic_H1_symplectic}
  Analyticity of the $a$- and $b$-loops and the symplectic $H_1$ basis on the elliptic curve. \emph{Cluster of 3 axioms.}
\end{lemma}

\section{Class 2d --- flagged (the deepest gaps)}

True-but-unproven: the canonical-differentials theorem. The route-D pipeline targets these.

\begin{lemma}[AX\_HyperellipticForm\_polynomial\_decomposition]
  \label{ax:AX-HyperellipticForm-polynomial-decomposition}
  \lean{Jacobians.Axioms.HyperellipticLiouville.AX_HyperellipticForm_polynomial_decomposition}
  Liouville L2: every form is $g(x)\,dx/y$, $\deg g < N/2-1$ (SP-7).
\end{lemma}

\begin{lemma}[AX\_HyperellipticOneForm\_eq\_form]
  \label{ax:AX-HyperellipticOneForm-eq-form}
  \lean{Jacobians.Axioms.HyperellipticLiouville.AX_HyperellipticOneForm_eq_form}
  Liouville L3: surjectivity onto low-degree forms (reduces to L2).
\end{lemma}

\chapter{Recently discharged}

Former axioms now proved in Lean --- the growing green frontier.

\begin{lemma}[Liouville L1: a holomorphic function on a compact complex manifold is locally constant]
  \label{done:liouville}
  \lean{Jacobians.Axioms.HyperellipticLiouville.liouville_compact_complex_manifold}
  \leanok
\end{lemma}

\begin{lemma}[growth $\Rightarrow$ polynomial (Liouville L2 step 4)]
  \label{done:growth}
  \lean{Jacobians.GeneralResults.differentiable_eq_polynomial_of_growth}
  \leanok
\end{lemma}

\begin{lemma}[odd-part difference quotient is analytic (route-D branch-point cancellation)]
  \label{done:oddpart}
  \lean{Jacobians.GeneralResults.analyticAt_dslope_oddPart}
  \leanok
\end{lemma}

